%% =====================================================================
%%  T-26-02  The Ordered Exit
%%  -------------------------------------------------------------------
%%  One idea per file. Core is mandatory. Slots in canonical order.
%%  Max 700 words. No layout commands. No filler. New source material
%%  for this entry goes in sources/B1/T-26-02.tex -- never by
%%  rewriting this file (docs/RULES.md R0.1).
%%  =====================================================================
%% @kind: topic
%% @id: T-26-02
%% @title: The Ordered Exit
%% @subtitle: Disengage before the structure turns on you
%% @chapter: c26
%% @order: 10
%% @status: stable
%% @sources: B1
%% @updated: 2026-09-19
%% @allow-rewrite: B1 ch.16 ingest 2026-09-19 -- committed scaffold replaced by the written entry (planned -> stable lifecycle)

\topic{T-26-02}{The Ordered Exit}{Disengage before the structure turns on you}

\begin{core}
Enter every engagement with its exit already designed. Take what the deal offers without taking the relationship apart, because grabbing for too much on one deal gets you cut out of the next one, and the habit you form with one counterpart is the habit you bring to all of them. Disengage while the structure --- your counterpart's willingness to deal again --- is still standing.
\end{core}
\begin{mechanism}
The source's cost chapter carries a pricing rule alongside its warnings: overreach on this deal and the counterpart will be sure to leave you out of the next one --- people catch on, and they end the relationship rather than renegotiate it. Why late exits fail is the chapter's harder finding: conduct does not change much from one relationship to the next. Treat one person as a tool and the treating spreads, because the trait, not the target, is the unit. The chapter's image is an eagle that spots a fish frozen in a block of ice drifting toward a waterfall, dives to snatch it, sticks, and goes over the falls still holding the ice. The snatch that was supposed to be the exception is the exception that ends the bird. Hence the exit must be ordered in advance: enough defined before the wanting starts, the number stopped at, the counterparty left whole --- not as virtue but as repeat business. An intact counterpart is a counterparty who deals again.
\end{mechanism}
\begin{conditions}
Strongest in repeat-trade worlds --- clients, colleagues, industries with memory --- and inside any institution where this quarter's counterpart is next quarter's gatekeeper. Weakens in genuine one-shots, which are rarer than they look; reputations quietly make every deal a repeat. The discipline buys nothing against a counterpart who is running a one-shot against you --- check their horizon before paying for a relationship they are not keeping.
\end{conditions}
\begin{application}
  \item Define the target before the engagement: what is enough, in numbers, dated.
  \item Write the exclusion first --- what you will not take, stated before the wanting starts.
  \item Stop at the number. The win is banked; the lap after the win is for the ego.
  \item When you feel the pull of just-a-little-more, treat it as the momentum alarm: that pull is the dive.
  \item Leave the door open by name --- say the next deal out loud before you leave this one.
\end{application}
\begin{feedback}
  \item Working: the counterpart initiates the second deal, or sends the referral unprompted.
  \item Working: terms tighten in your favour over time --- the market has learned you stop at the mark.
  \item Failing: you got the terms, and the silence that followed them. The account was closed at the table.
  \item Failing: your conduct in unrelated relationships drifts toward the same take. The eagle is already in the ice.
\end{feedback}
\begin{failure}
The discipline reads as leaving money on the table, and sharper counterparts will fill the gap you leave --- ordered exits are a discount paid for franchise, and someone always bids against discounts. Half-ordered exits are the worst outcome on the board: a person who takes a little too much has paid the reputation cost without buying the return. And the momentum law means the discipline cannot be improvised at the end --- the exit that was not designed at entry is not an exit, it is a slide.
\end{failure}
\begin{countermeasures}
  \item Price a counterpart who over-reaches once as data; twice as character. The second overreach is the slide made visible.
  \item Refuse to be the ice: structure your own exposure so a counterpart's dive cannot be funded out of your side of the deal.
  \item In your own shop, measure deals by whether the counterpart would redo them --- that metric, not gross extraction, predicts the franchise.
  \item When you catch yourself renegotiating an agreed exit with yourself, say the original number out loud. The slide hates an audience (T-25-01).
\end{countermeasures}

\sources{B1}
\seesources
\seealso{T-25-01, T-27-01, T-16-03}
